% =====================================================================
%  Methodology -- MeshGraphNets-V and cHI-MGNflow
%  Compile through main.tex.
% =====================================================================

\section{Methodology}
\label{sec:method}

\subsection{Problem statement: predicting a conditional law}
\label{sec:method:problem}

Let $g$ denote everything a learned simulator observes about a scene: the mesh
$(\mathcal{V},\mathcal{E})$, the nodal coordinates, the boundary-condition and
material channels, and any scalar conditions supplied alongside them. Let
$\mathbf{y}\in\mathbb{R}^{N\times F}$ be the field to be predicted at the $N$
mesh nodes. The deterministic simulator lineage---GNS, MeshGraphNets, and the
hierarchical variants including HI-MGN---learns a map $\hat{\mathbf{y}} =
f_\theta(g)$ under a pointwise regression loss,
\begin{equation}
  \mathcal{L}_{\mathrm{MSE}}(\theta)
  = \mathbb{E}_{(g,\mathbf{y})}\bigl\|f_\theta(g)-\mathbf{y}\bigr\|_2^2 ,
  \label{eq:mse}
\end{equation}
whose population minimiser is the conditional mean
$f^\star(g)=\mathbb{E}[\mathbf{y}\mid g]$.

The conditional mean is optimal for squared-error point prediction, but does
not specify uncertainty when $p(\mathbf{y}\mid g)$ is broad. Write the data-generating
process as
\begin{equation}
  \mathbf{y} = \mathcal{S}(g,\boldsymbol{\xi}),
  \qquad \boldsymbol{\xi}\sim p(\boldsymbol{\xi}\mid g),
  \label{eq:generative}
\end{equation}
where $\mathcal{S}$ is the physical solver and $\boldsymbol{\xi}$ is a latent
variable that governs the response but is \emph{not} part of $g$: an
unmeasured geometric imperfection, a process-to-process manufacturing
deviation, a microstructural realisation. The quantity of engineering interest
is then the conditional law
\begin{equation}
  p(\mathbf{y}\mid g)
  = \int p\bigl(\mathbf{y}\mid g,\boldsymbol{\xi}\bigr)\,
         p(\boldsymbol{\xi}\mid g)\,\mathrm{d}\boldsymbol{\xi},
  \label{eq:marginal}
\end{equation}
and a model trained under \eqref{eq:mse} returns a single point estimate.
When \eqref{eq:marginal} is multimodal, as can occur in buckling mode
selection, the conditional mean need not lie in its support: averaging
physically realisable states can produce a field that is itself unrealizable. Tolerance
analysis, yield prediction and reliability assessment all require the width and
shape of \eqref{eq:marginal}, not its centroid.

We therefore ask a model to emit \emph{samples} and compare the generated and
reference ensembles. Section~\ref{sec:protocol} distinguishes the historical
SAOI scalar-distance metric from the implemented buckling energy score. The
two mesh-native constructions below place stochasticity in a graph latent or
directly in the field. This section describes the current implementation;
the historical SAOI results in Section~\ref{sec:results} are not a validation
of every subsequently added option.

\subsection{Shared backbone}
\label{sec:method:backbone}

Both models inherit the encode--process--decode structure of MeshGraphNets and
the hierarchical processor of HI-MGN. Node and edge features are embedded by
MLPs into latent vectors $\mathbf{v}_i$ and $\mathbf{e}_{ij}$, and each message
passing (MP) block performs
\begin{align}
  \mathbf{e}_{ij}' &= f_E\bigl(\mathbf{e}_{ij},\,\mathbf{v}_i,\,\mathbf{v}_j\bigr),
  \label{eq:mp-edge}\\
  \mathbf{v}_i'    &= f_V\Bigl(\mathbf{v}_i,\;
                        \textstyle\sum_{j\in\mathcal{N}(i)}\mathbf{e}_{ij}'\Bigr),
  \label{eq:mp-node}
\end{align}
with residual connections, after which a decoder MLP reads the nodal output.
The receptive field expands by at most one mesh hop per block; communicating
between arbitrary nodes through this path alone requires depth at least the
graph diameter. HI-MGN replaces the flat stack with a multiscale processor.
Farthest-point sampling chooses seeds and multi-source breadth-first search
assigns nodes to their nearest seed in graph distance. Coarse edges are
induced by fine edges crossing cluster boundaries. In
\texttt{voronoi\_seedmean} mode, features are mean-pooled while coarse
positions use the seed nodes. MP blocks run at each level, and a learned
graph interpolation network reconstructs fine-resolution features on the way
back up. For example, the current SAOI arm-1 configurations specify two
coarsening levels with $1000$ and $100$ requested coarse nodes, hidden width
$128$, and a $4,6,8,6,4$ MP schedule. Other arms vary capacity or conditioning.
Sharing this hierarchy does not make the full models identical: posterior,
prior, input dimensions, modulation and training objectives also differ.

A note on the input contract. The \texttt{cond\_var} channels following the
state rows in \texttt{nodal\_data} are
\emph{conditions}: quantities the model reads but never predicts (process
settings, known boundary values). They enter $\mathbf{x}$ alongside the state,
positional and node-type features, receive their own normalisation statistics,
and are retained as known inputs during an autoregressive rollout. For the
static $T=1$ contract, the state input is zero and the target is the absolute
output field; temporal data instead use the next-step increment as target.
The conditions are part of $g$; the latent $\boldsymbol{\xi}$ of
\eqref{eq:generative} is, by construction, not.

\subsection{MeshGraphNets-V: a latent-variable simulator}
\label{sec:method:mgnv}

MeshGraphNets-V (MGN-V) uses a conditional variational-encoder architecture
with an aggregate-posterior MMD objective, rather than a KL-based ELBO. A graph-level latent
$\mathbf{z}\in\mathbb{R}^{K\times d}$ carries the sample-to-sample spread while
the processor carries what geometry and boundary conditions determine. The
latent is organised as $K=L+1$ \emph{slots}, one per hierarchy level plus the
fine level, each of width $d$; the current SAOI configuration uses
$K=3$, $d=16$, i.e.\ $48$ latent coordinates in total. These are graph-level
slots injected at different levels, not a latent vector at every coarse node.

\paragraph{Posterior.}
A GNN encoder sees the normalised target and edge features and, when
\texttt{vae\_graph\_aware=True} as in the current SAOI configurations, also
the normalised input-node features. It runs message passing and
pools to a single graph vector by attention, from which all slots are emitted:
\begin{equation}
  q_\phi(\mathbf{z}\mid \mathbf{y},g)
  = \mathcal{N}\bigl(\boldsymbol{\mu}_\phi(\mathbf{y},g),\,
                     \operatorname{diag}\boldsymbol{\sigma}^2_\phi(\mathbf{y},g)\bigr).
  \label{eq:posterior}
\end{equation}

\paragraph{Decoder.}
The HI-MGN simulator itself is the decoder. Each hierarchy level receives its
corresponding slot at every MP block. The implementation supports a linear
concatenation fuser and AdaLN-Zero modulation; the current SAOI configurations
select AdaLN-Zero. The posterior variance is diagonal conditional on its
inputs, although mixing over targets can induce dependence between slots.

\paragraph{Conditional prior.}
At inference the target is unavailable. With the learned conditional prior
enabled, $\mathbf{z}$ is drawn from $p_\psi(\mathbf{z}\mid g)$. Its default family is conditional
flow matching in latent space: a velocity network transports
$\mathcal{N}(\mathbf{0},\mathbf{I})$ onto the posterior cloud for that graph,
conditioned on a pooled graph embedding, and is integrated by a Heun or Euler
ODE solver. All slots are flattened jointly, so cross-slot dependence is
learned rather than assumed factorised. The implemented path retains endpoint
noise with standard deviation $10^{-4}$ in the prior's working coordinates.
An optional frozen-prior training phase freezes the encoder and decoder,
standardises posterior latents using training-set moments, and fits only the
prior; sampling undoes that affine standardisation. A legacy
graph-conditioned Gaussian mixture and standard-normal latent sampling are
also available. They are distinct inference choices, not equivalent priors.

\paragraph{Objective.}
For joint training with the flow-matching prior, the current loss is
\begin{equation}
  \mathcal{L}_{\text{MGN-V}}
  = \alpha_{\mathrm{recon}}\,\mathcal{L}_{\mathrm{recon}}
  + \lambda_{\mathrm{MMD}}\,\mathcal{L}_{\mathrm{MMD}}
  + \beta_{\mathrm{aux}}\,\mathcal{L}_{\mathrm{aux}}
  + w_{\mathrm{prior}}\,\mathcal{L}_{\mathrm{FM}},
  \label{eq:mgnv-loss}
\end{equation}
where $\mathcal{L}_{\mathrm{recon}}$ is the channel-weighted, node-averaged
posterior reconstruction error in normalised target coordinates (MSE in the
current SAOI configurations; Huber is also supported).
$\mathcal{L}_{\mathrm{MMD}}$ averages a multi-bandwidth RBF MMD squared over
slots, comparing minibatch posterior samples with independent
$\mathcal{N}(0,I)$ samples. It uses the biased V-statistic, including diagonal
kernel terms, with fixed or detached median-based bandwidths. This is
aggregate-posterior regularisation in the InfoVAE sense
\citep{zhao2017infovae}; it does not match each conditional posterior to the
learned $p_\psi(\mathbf{z}\mid g)$, or constrain the joint cross-slot law.
$\mathcal{L}_{\mathrm{aux}}$ is the graph-averaged squared error between
predicted and target peak-to-valley values of a selected normalised output
channel. For $T=1$ this is the field statistic; for temporal data it is the
increment statistic. It encourages reconstruction of that statistic without
guaranteeing calibrated spread or a mutual-information lower bound.
$\mathcal{L}_{\mathrm{FM}}$ is the flow-matching
loss that fits the conditional prior to a \emph{fresh} posterior sample drawn
each step. A switch $g_{\mathrm{enc}}\in[0,1]$ scales the gradient that
$\mathcal{L}_{\mathrm{FM}}$ sends back through the encoder; at
$g_{\mathrm{enc}}=0$ the prior chases a detached posterior, and at
$g_{\mathrm{enc}}=1$ the encoder also receives the velocity-regression
gradient through the path point and target.
Section~\ref{sec:method:shrink} shows why that second coupling is not the
variational rate term it resembles.

\subsection{cHI-MGNflow: flow matching in field space}
\label{sec:method:flow}

cHI-MGNflow places the stochasticity at the output instead of in a bottleneck.
It learns a conditional flow directly on the $N\times F$ field. Let
$\mathbf{y}_1$ be the normalised target and
$\mathbf{y}_0\sim\mathcal{N}(\mathbf{0},\mathbf{I})$ an independent Gaussian
source drawn at every node and channel. With $s=1-\sigma_{\min}$ the
conditional path and its velocity are, with the implemented
$\sigma_{\min}=10^{-4}$,
\begin{align}
  \mathbf{y}_t &= (1-s\,t)\,\mathbf{y}_0 + t\,\mathbf{y}_1,
  \label{eq:flow-path}\\
  \mathbf{u}   &= \mathbf{y}_1 - s\,\mathbf{y}_0 ,
  \label{eq:flow-target}
\end{align}
and the network regresses the velocity,
\begin{equation}
  \widehat{\mathcal{L}}_{\mathrm{flow}}
  = \frac{\sum_b w(t_b)\sum_{i=1}^{N_b}\sum_f a_f
    \bigl(v_\theta(\mathbf{y}_{t_b},t_b,g_b)_{if}-u_{bif}\bigr)^2}
    {\sum_b N_b w(t_b)} ,
  \label{eq:flow-loss}
\end{equation}
where $a_f$ are nonnegative channel weights normalised to sum to one, and
$b$ indexes graphs in a collated training minibatch. The implementation normalises by the
realised sum of node-time weights (with a numerical denominator floor), not
by an expectation over $t$. Without time weighting, $w(t)=1$. The flow time
$t$ is drawn \emph{per graph}, not per node: the velocity field
is a function of the whole field state, so mixing times inside one graph would
depart from the stated whole-field path. The hierarchical
processor is traversed once per velocity evaluation, with $t$ entering through
AdaLN modulation and the noisy field entering alongside the conditions.

The default \texttt{flow\_head=v} predicts velocity directly. An optional
data-prediction weighting, called \texttt{x0} in the configuration, uses
$w(t)=(1-st)^2$. The per-node loss identity is exact:
\begin{equation}
  \hat{\mathbf{y}}_1 = s\,\mathbf{y}_t + (1-s\,t)\,\hat{\mathbf{u}}
  \quad\Longrightarrow\quad
  \bigl\|\hat{\mathbf{y}}_1-\mathbf{y}_1\bigr\|^2
  = (1-s\,t)^2\bigl\|\hat{\mathbf{u}}-\mathbf{u}\bigr\|^2 ,
  \label{eq:x0-equiv}
\end{equation}
The batch normalisation in \eqref{eq:flow-loss} still applies. Separately,
the current code supports \texttt{flow\_head=x}, whose network output is a
clean-field estimate $h_\theta$ and whose returned velocity is
\begin{equation}
 v_\theta = \frac{h_\theta-s\mathbf{y}_t}
 {\max(1-st,\varepsilon_{\mathrm{head}})} .
 \label{eq:flow-x-head}
\end{equation}
The default floor is $\varepsilon_{\mathrm{head}}=0.05$. This conversion
equals the unfloored path identity only where $1-st\geq
\varepsilon_{\mathrm{head}}$. Without this floor, the gain reaches $10^4$ at
$t=1$; it is finite at the implemented positive $\sigma_{\min}$ and becomes
singular only in the zero-noise limit. Thus output-head choice and time-loss
weighting are separate options.

Time sampling can be uniform or logit-normal, with an optional probability
of pinning a graph to $t=0$. At that endpoint, independence of
$\mathbf{y}_0$ and $\mathbf{y}_1$ gives the population regression optimum
$v^*(\mathbf{y}_0,0,g)=\mathbb{E}[\mathbf{y}_1\mid g]-s\mathbf{y}_0$.
The implemented one-evaluation readout is
$\mathbf{y}_0+v_\theta(\mathbf{y}_0,0,g)$, which at this optimum equals
$\mathbb{E}[\mathbf{y}_1\mid g]+\sigma_{\min}\mathbf{y}_0$. Its error also
depends on how well the trained network fits $t=0$.

Sampling integrates the learned ODE from $t=0$ to $t=1$ with Euler or Heun,
holding the graph hierarchy fixed. Even with the exact marginal velocity
and exact integration, the stated path ends at
$\mathbf{y}_1+\sigma_{\min}\mathbf{y}_0$, so its target law is the
Gaussian-smoothed data law. A finite network and finite integration budget
introduce additional approximation error.

\subsection{A contraction mechanism in trainable-target flow matching}
\label{sec:method:shrink}

A tempting reading of \eqref{eq:mgnv-loss} is that opening the encoder gradient
($g_{\mathrm{enc}}=1$) turns $\mathcal{L}_{\mathrm{FM}}$ into the
$D_{\mathrm{KL}}\bigl(q(\mathbf{z}\mid\mathbf{y},g)\,\|\,p(\mathbf{z}\mid
g)\bigr)$ rate term of a conditional VAE, making \eqref{eq:mgnv-loss} an ELBO.
It does not. A Gaussian counterexample exhibits the difference.

The flow-matching theorem \citep{lipman2023flowmatching} equates the gradients of the conditional and marginal
velocity-regression objectives \emph{with respect to the velocity model's
parameters, for a fixed target distribution}. It says nothing about gradients
with respect to parameters that \emph{move} the target. A VAE rate term
contains the posterior entropy and vanishes whenever $q=p$. The optimised
flow-matching MSE need not vanish when its endpoint distribution is matched.

The consequence is quantitative. Consider the scalar Gaussian instance
$Z_0\sim\mathcal{N}(0,1)$, $Z_1\sim\mathcal{N}(0,a^2)$ independent, with the
linear path $Z_t=(1-t)Z_0+tZ_1$ and velocity target $U=Z_1-Z_0$. Here $a$
plays the role of the trainable posterior scale, with $a>0$. This example
sets $\sigma_{\min}=0$ to give a simple closed form. Since $(U,Z_t)$ is jointly
Gaussian, the optimal regressor is the conditional expectation and its
irreducible loss is the conditional variance,
\begin{equation}
  \operatorname{Var}\bigl(U \mid Z_t\bigr)
  = (1+a^2) - \frac{\bigl(a^2t-(1-t)\bigr)^2}{(1-t)^2+a^2t^2}
  = \frac{a^2}{(1-t)^2+a^2t^2}.
  \label{eq:cond-var}
\end{equation}
Integrating over a uniform time schedule gives the loss floor attained by a
\emph{perfectly} trained velocity model,
\begin{equation}
  \int_0^1 \operatorname{Var}\bigl(U\mid Z_t\bigr)\,\mathrm{d}t
  = \int_0^1 \frac{a^2}{(1-t)^2+a^2t^2}\,\mathrm{d}t
  = \frac{\pi a}{2}.
  \label{eq:floor}
\end{equation}
The floor is \emph{linear in $a$}. Meanwhile a prior that matches the target
exactly has $D_{\mathrm{KL}}=0$ for every $a$. So the optimised flow-matching
objective can always be decreased by shrinking the target distribution, even
though the density-matching problem it is supposed to represent is already
solved perfectly at every scale. In this Gaussian example, allowing the
encoder scale to optimise the best-achievable FM loss therefore creates
contraction pressure even with unlimited prior capacity. For the implemented
positive endpoint noise, the corresponding variance is
$a^2/[(1-st)^2+a^2t^2]$; its derivative with respect to $a$ is
$2a(1-st)^2/[(1-st)^2+a^2t^2]^2>0$, so the scale incentive remains in this
example. This calculation concerns the FM term alone, without reconstruction,
MMD, variance floors or latent standardisation.

The counterexample shows that a trainable-target FM loss is not generally a
VAE rate term. It does not establish the net effect of the coupling in a
nonlinear, jointly trained model, nor rule out an empirical benefit. The
historical SAOI note reports a small observed contrast for this factor
(Section~\ref{sec:results}); the incomplete, unreplicated sweep cannot prove
that the causal effect is zero or identify the cause of under-dispersion.
